% Bagian: incompleteness
% Bab: theories-computability
% Subbagian: inseparability

\documentclass[../../../include/open-logic-section]{subfiles}

\begin{document}

\olfileid[id]{inc}{tcp}{ins}

\olsection{Kalimat yang Dapat Dibuktikan dan Disangkal dalam $\Th{Q}$ Tak
  Terpisahkan secara Komputabel}


Misalkan $\Th{\bar Q}$ sebagai himpunan kalimat yang {\em negasinya} dapat
dibuktikan dalam~$\Th{Q}$, yakni $\Th{\bar Q} = \Setabs{!A}{\Th{Q} \Proves
  \lnot !A}$. Ingat bahwa himpunan saling lepas~$A$ dan~$B$ disebut tak
terpisahkan secara komputabel jika tidak ada himpunan dapat dihitung~$C$
sedemikian sehingga $A \subseteq C$ dan $B \subseteq \Complement{C}$.

\begin{lem}
$\Th{Q}$ dan $\Th{\bar Q}$ tak terpisahkan secara komputabel.
\end{lem}

\begin{proof}
Andaikan $C$ suatu himpunan dapat dihitung sedemikian sehingga $\Th{Q}
\subseteq C$ dan $\Th{\bar Q} \subseteq \Complement{C}$. Misalkan $R(x,y)$
sebagai relasi
\begin{quote}
$x$ mengodekan suatu formula~$!D(u)$ dan $!D(\num y)$ berada dalam~$C$.
\end{quote}
Kita akan menunjukkan bahwa $R(x,y)$ merupakan relasi dapat dihitung
universal, yang menghasilkan kontradiksi.

Andaikan $S(y)$ dapat dihitung dan direpresentasikan oleh~$!D_S(u)$
dalam~$\Th{Q}$. Maka
\begin{eqnarray*}
S(\num n) & \lif & \Th{Q} \Proves !D_S(\num n) \\
& \lif & !D_S(\num n) \in C
\end{eqnarray*}
dan
\begin{eqnarray*}
\lnot S(\num n) & \lif & \Th{Q} \Proves \lnot !D_S(\num n) \\
& \lif & !D_S(\num n) \in \Th{\bar Q} \\
& \lif & !D_S(\num n) \not\in C
\end{eqnarray*}
Jadi $S(y)$ ekuivalen dengan $R(\#(!D_S(\num u)),y)$.
\end{proof}

\end{document}
